\documentclass{article}
\usepackage[T1]{fontenc}
\usepackage{lmodern}
\usepackage{graphicx}
\usepackage{amsmath,amssymb}
\usepackage[a4paper,margin=2cm]{geometry}
\usepackage[absolute,overlay]{textpos}
\setlength{\TPHorizModule}{1mm}
\setlength{\TPVertModule}{1mm}
\usepackage[indonesian]{babel}
\usepackage{hyperref}
\setlength{\emergencystretch}{2em}
% unit: Halaman 15

\begin{document}

% === Halaman 15 (llm) ===

Alice melakukan enkripsi dengan kunci public Bob ($y = 1185, g = 2, p = 2357$):
\[ a = g^k \pmod p = 2^{1520} \pmod{2357} = 1430 \]
\[ b = y^k m \pmod p = 1185^{1520} \cdot 2035 \pmod{2357} = 697 \]

Jadi, cipherteks yang dihasilkan adalah $(1430, 697)$.\nAlice mengirim cipherteks ini kepada Bob.

\textbf{(c) Dekripsi (Oleh Bob)}

Bob melakukan dekripsi dengan kunci privatnya ($x = 1751, p = 2357$):
\[ (a^x)^{-1} = a^{p-1-x} \pmod p = 1430^{2357-1-1751} \pmod{2357} = 1430^{605} \pmod{2357} = 872 \]
\[ m = b/a^x \pmod p = b \cdot (a^x)^{-1} \pmod p = 697 \cdot 872 \pmod{2357} = 2035 \]

Bob mendapatkan kembali plainteks $m = 2035$ yang dikirim oleh Alice.

\end{document}
